% !PRINTED 279
\weberpage{279}
Now on the other hand $\omega$ can again be taken so small that the roots of $L_n' = 0$ differ as little as we please from those of $L_n = 0$ (\S.\,40).

A double root of $L_n=0$ may indeed pass over into two simple roots of $L_n'=0$; but since $L_n'=0$ has no imaginary roots, these are real; and the same holds when $L_n=0$ has multiple roots.

\gesperrt{Accordingly the theorems I, II retain their validity when the assumption is given up that in the chain of the $L_k$ no consecutive terms vanish; only the multiple roots must then be counted according to their multiplicity.}

\begin{center}
\S.\,89.

\gesperrt{The Sturm functions.}
\end{center}

After these particular examples we turn to the consideration of the procedure by which one obtains a Sturm chain in all cases\footnote{\gesperrt{Sturm}, M\'em. sur la r\'esolution des \'equations num\'eriques. M\'em. de l'acad\'emie de Paris. Sav. \'etrang. VI, 1835. Extract in Bull. de Ferussac XI, 1829.}. The manner of forming these functions is in principle extraordinarily simple, even though in practical execution it is usually not feasible.

We restrict ourselves here to the consideration of equations \gesperrt{without multiple roots}, or we assume that before the application $f(x)$ has been freed of every common factor with its derivative $f'(x)$.

If we then set
\weq{1}{f_1(x) = f'(x)}
the condition \S.\,86, 4. is certainly satisfied. We now proceed as if it were a matter of finding the greatest common divisor of $f(x)$ and $f_1(x)$, reversing each time the sign of the remainder; we thus form by division the equations
\weq{2}{\begin{gathered}
f = q_1 f_1 - f_2\\
f_1 = q_2 f_2 - f_3\\
\cdots\cdots\\
f_{m-2} = q_{m-1} f_{m-1} - f_m,
\end{gathered}}
%CONT>>
